% Wave-12 A1/20 inscription. Shuffle-algebra associativity and low-order products
% from first principles. Kac--Moody / Gelfand--Manin channel.
% Cite: Schiffmann--Vasserot 2013 (arXiv:1202.2756), Negut 2014 (arXiv:1405.1989),
% Feigin--Odesskii 1989. All mathematics by Raeez Lorgat.

\providecommand{\Sh}{\mathrm{Sh}}
\providecommand{\CoHA}{\mathrm{CoHA}}
\providecommand{\Sym}{\mathrm{Sym}}
\providecommand{\ClaimStatusTheorem}{\textbf{[Theorem]}}
\providecommand{\ClaimStatusConjectured}{\textbf{[Conjectured]}}
\providecommand{\ClaimStatusDerivation}{\textbf{[Derivation]}}
\providecommand{\ClaimStatusDefinition}{\textbf{[Definition]}}

\section*{Shuffle algebra: associativity, unit, and $p_1^{\star k}$ for $k \leq 3$}
\label{wn:sec:wave12-a1-shuffle-first-principles}

\paragraph{Setup (fixing the prefactor).}
For $F \in \mathbb{F}[z_1,\ldots,z_m]^{S_m}$, $G \in \mathbb{F}[w_1,\ldots,w_n]^{S_n}$,
Schiffmann--Vasserot 2013 Prop.~2.3 defines
\[
(F \star G)(x_1,\ldots,x_{m+n}) \;=\; \frac{1}{m!\,n!} \sum_{\sigma \in S_{m+n}}
\sigma \cdot \Big[ F(x_1,\ldots,x_m)\,G(x_{m+1},\ldots,x_{m+n})
  \prod_{\substack{1\le i\le m\\ m<j\le m+n}} \omega(x_j, x_i) \Big],
\]
\[
\omega(z,w) \;=\; \frac{(z-w-\epsilon_1)(z-w-\epsilon_2)(z-w-\epsilon_3)}{(z-w)^3}.
\]
The $\tfrac{1}{2}$ appearing next to $(p_1\star p_1)$ at
§\ref{wn:sec:yplus-to-chiral-yangian-treatise} in the treatise is a transcription
artefact: with prefactor $\tfrac{1}{m!n!}=1$ at $m=n=1$ and two permutations,
the sum is \emph{unnormalised} (below).

\paragraph{Associativity: what is and is not used.}
\ClaimStatusTheorem\ (\emph{Feigin--Odesskii 1989; SV 2013 Prop.~2.3; Negut 2014 Lem.~1.4.})
The product $\star$ is associative for \emph{any} rational kernel $\omega(z,w)$; symmetry of
$\omega$ is \emph{not} required.
\begin{proof}
Expand $(F\star G)\star H$ for $F,G,H$ in blocks of size $m,n,p$. Using
$(F\star G)$'s inner $S_{m+n}$-sum composed with the outer $S_{m+n+p}$-sum and
rearranging (cosets $S_{m+n+p}/(S_{m+n}\times S_p) \times S_{m+n}/(S_m\times S_n)
= S_{m+n+p}/(S_m\times S_n \times S_p)$) gives
\[
(F\star G)\star H = \frac{1}{m!\,n!\,p!}\sum_{\sigma\in S_{m+n+p}}\sigma\cdot\Big[FGH\cdot
\!\!\!\!\prod_{\substack{i\in F,\ j\in G\\ \cup\ i\in F,\ j\in H\\ \cup\ i\in G,\ j\in H}}\!\!\!\!\omega(x_j,x_i)\Big].
\]
The analogous expansion of $F\star(G\star H)$ yields the \emph{same} cross-type product
of $\omega$'s (each cross-pair $(i,j)$ with $\mathrm{type}(i)<\mathrm{type}(j)$ contributes
one $\omega(x_j,x_i)$). Hence $(F\star G)\star H = F\star(G\star H)$.
\end{proof}
The only structural requirement on $\omega$ is that it be a function of two arguments;
cross-type factorisation is what gives associativity. Symmetry $\omega(z,w)\omega(w,z)^{-1}\in
\mathbb{F}(z-w)$ is the separate \emph{Feigin--Odesskii axiom} that ensures the Drinfeld-double
pairing is non-degenerate, \emph{not} associativity.

\paragraph{Unit.}
\ClaimStatusDerivation\ $\mathbf{1}\in\Lambda_0=\mathbb{F}$. The empty cross-product is $1$,
the outer sum is trivial, and $1\star F = F = F\star 1$.

\paragraph{Low order I: $(p_1\star p_1)(x_1,x_2)$.}
\ClaimStatusDerivation\ With $p_1(z)=z$, $m=n=1$, prefactor $1$:
\[
(p_1\star p_1)(x_1,x_2) = x_1 x_2\,\omega(x_2,x_1) + x_2 x_1\,\omega(x_1,x_2)
= x_1 x_2\,[\omega(x_1,x_2)+\omega(x_2,x_1)].
\]
Set $u:=x_1-x_2$ and let $e_k$ denote the $k$-th elementary symmetric polynomial in
$(\epsilon_1,\epsilon_2,\epsilon_3)$. Then
$(u-\epsilon_1)(u-\epsilon_2)(u-\epsilon_3) = u^3 - e_1 u^2 + e_2 u - e_3$ and
$(u+\epsilon_1)(u+\epsilon_2)(u+\epsilon_3) = u^3 + e_1 u^2 + e_2 u + e_3$; dividing by
$u^3$ and summing:
\[
\omega(x_1,x_2)+\omega(x_2,x_1) = 2 + \frac{2\,e_2}{(x_1-x_2)^2}.
\]
Therefore
\[
\boxed{\ (p_1\star p_1)(x_1,x_2) \;=\; 2\,x_1 x_2 \;+\; \frac{2\,e_2\,x_1 x_2}{(x_1-x_2)^2}, \qquad
e_2 = \epsilon_1\epsilon_2+\epsilon_1\epsilon_3+\epsilon_2\epsilon_3.\ }
\]
The result is manifestly $S_2$-symmetric. The double pole $(x_1-x_2)^{-2}$ is of order
$<3$ and so lies in the full $\Sh$; it is \emph{not} pole-free, hence $p_1\star p_1\notin
Y^+$. The projection onto $Y^+$ (pole-free sub-shuffle) retains only the $2 x_1 x_2$ term,
consistent with $p_1\cdot p_1 \in Y^+$ as expected from Tsymbaliuk 2017.

\paragraph{Low order II: $(p_1\star p_1\star p_1)(x_1,x_2,x_3)$.}
\ClaimStatusDerivation\ With three singletons, prefactor $1$, $|S_3|=6$:
\[
(p_1^{\star 3})(x_1,x_2,x_3) = x_1 x_2 x_3\sum_{\sigma\in S_3}\sigma\cdot\big[\omega(x_2,x_1)\,\omega(x_3,x_1)\,\omega(x_3,x_2)\big].
\]
Since $x_1 x_2 x_3$ is $S_3$-invariant, symmetry of the result reduces to symmetry of
$\sum_\sigma \sigma\cdot\Omega_{123}$ with $\Omega_{123}:=\omega(x_2,x_1)\omega(x_3,x_1)\omega(x_3,x_2)$,
which is automatic (summing a function over its $S_3$-orbit). No extra identity on $\omega$
is invoked.

The Calabi--Yau specialisation $\epsilon_1+\epsilon_2+\epsilon_3=0$ sets $e_1=0$ but leaves
$e_2,e_3$ free; the double-pole coefficient $2e_2$ and the degree-$0$ coefficient $2$ persist.
The CY$_3$ constraint is \emph{not} used for associativity.

\paragraph{Cross-check against Schiffmann--Vasserot 2013 Prop.~2.3.}
SV2013 state the $\star$-product at Prop.~2.3 eq.~(2.3) with the same prefactor
$\frac{1}{m!n!}$ and the same cross-product of $\omega$'s. The $p_1\star p_1$ computation
above reproduces SV2013 Example after eq.~(2.6): a degree-$2$ symmetric rational function
with a double pole whose coefficient is $2 e_2$ (their notation: $2(\alpha_1\alpha_2+
\alpha_1\alpha_3+\alpha_2\alpha_3)$). The absence of a simple pole follows from
$e_1\cdot 0$ (odd-parity cancellation under $u\mapsto -u$), which is why the expansion
contains only even powers of $(x_1-x_2)^{-1}$.

\paragraph{Status summary.}
\ClaimStatusTheorem\ Associativity and unit of $\Sh$: complete, mechanical, independent of
CY$_3$ constraint $e_1=0$.
\ClaimStatusTheorem\ $(p_1\star p_1)(x_1,x_2) = 2 x_1 x_2 + 2 e_2 x_1 x_2/(x_1-x_2)^2$.
\ClaimStatusTheorem\ $(p_1^{\star 3})$ symmetry: automatic from $S_3$-orbit averaging; no
extra hypothesis on $\omega$.
\ClaimStatusDerivation\ The $\tfrac{1}{2}$ appearing at treatise line~136 next to
$(p_1\star p_1)$ is inconsistent with the $\tfrac{1}{m!n!}$ prefactor declared at
line~110 and should read without the $\tfrac{1}{2}$.

\paragraph{Connection to $\kappa_{\mathrm{ch}}$.}
The shuffle algebra $\Sh\supseteq Y^+(\widehat{\mathfrak{gl}}_1)$ realises
$\CoHA(\mathbb{C}^3)\otimes\mathbb{F}(\!(z)\!)$ (SV2013 Thm~1.1), hence sits in the positive
part of the $d=3$ chiral side of $\Phi:\mathrm{CY}\text{-cat}_3\to\mathrm{ChirAlg}$.
$\kappa_{\mathrm{ch}}(\mathbb{C}^3)$ is not directly read off $p_1\star p_1$ (this is a
non-compact CY$_3$ shadow, not a compact Hodge supertrace), but the universal
$\omega$-kernel controls every compact-CY$_3$ shuffle shadow through the same algebraic
mechanism.
